\documentclass[11pt,oneside,headings=big]{scrreprt}

% Shared visual style for the repository's LaTeX documents.
% Keep document-specific packages, metadata, and content in each source file.

\usepackage[letterpaper,margin=0.86in,headheight=16pt,footskip=28pt]{geometry}
\usepackage[T1]{fontenc}
\usepackage{newtxtext,newtxmath}
\usepackage{microtype}
\usepackage{graphicx}
\usepackage{booktabs,tabularx,longtable,array}
\usepackage[table]{xcolor}
\usepackage[most]{tcolorbox}
\usepackage{enumitem}
\usepackage{scrlayer-scrpage}
\usepackage{caption}
\usepackage{subcaption}
\usepackage{hyperref}

\definecolor{ImuNavy}{HTML}{17324D}
\definecolor{ImuBlue}{HTML}{2E6F9E}
\definecolor{ImuTeal}{HTML}{168A8A}
\definecolor{ImuOrange}{HTML}{D9792B}
\definecolor{ImuGold}{HTML}{D9A441}
\definecolor{ImuInk}{HTML}{1E252B}
\definecolor{ImuSlate}{HTML}{52606D}
\definecolor{ImuPale}{HTML}{EEF4F7}
\definecolor{ImuPaleTeal}{HTML}{EAF6F5}
\definecolor{ImuPaleGold}{HTML}{FFF7E6}
\definecolor{ImuRule}{HTML}{CBD7DE}
\definecolor{ImuCode}{HTML}{F5F7F8}

\setkomafont{disposition}{\color{ImuNavy}\bfseries}
\setkomafont{chapter}{\color{ImuNavy}\bfseries\LARGE}
\setkomafont{section}{\color{ImuBlue}\bfseries\Large}
\setkomafont{subsection}{\color{ImuTeal}\bfseries\large}
\setkomafont{caption}{\small}
\setkomafont{captionlabel}{\color{ImuNavy}\bfseries}
\renewcommand*{\chapterformat}{%
    \colorbox{ImuNavy}{\parbox[c][1.65em][c]{2.35em}{\centering\color{white}\thechapter}}\enskip
}

\clearpairofpagestyles
\ihead{\textsc{IMU Error Model}}
\ohead{\headmark}
\cfoot*{\pagemark}
\addtokomafont{pageheadfoot}{\small\color{ImuSlate}}
\automark[section]{chapter}

\setlist[itemize]{leftmargin=1.4em,itemsep=0.25em,topsep=0.35em}
\setlist[enumerate]{leftmargin=1.6em,itemsep=0.25em,topsep=0.35em}
\renewcommand{\arraystretch}{1.18}
\setlength{\tabcolsep}{3pt}
\emergencystretch=1.5em

\newcolumntype{Y}{>{\raggedright\arraybackslash}X}
\newcolumntype{L}[1]{>{\raggedright\arraybackslash}p{#1}}
\newcolumntype{C}[1]{>{\centering\arraybackslash}p{#1}}
\newcommand{\software}[1]{\texttt{#1}}


\hypersetup{
    colorlinks=true,
    linkcolor=ImuBlue,
    urlcolor=ImuOrange,
    pdftitle={IMU Error Model Profile Showcase},
    pdfauthor={imu-error-model project},
    pdfsubject={Example IMU profiles and generated analysis plots}
}

\graphicspath{{../artifacts/showcase/}}

\ihead{\textsc{IMU Profile Showcase}}

\newtcolorbox{notionalbox}[1][]{
    enhanced,breakable,
    colback=ImuPaleGold,colframe=ImuGold,
    boxrule=0.8pt,arc=2pt,
    left=8pt,right=8pt,top=7pt,bottom=7pt,
    title={Interpretation boundary},
    fonttitle=\bfseries\color{ImuNavy},
    #1
}
\newtcolorbox{methodbox}[1][]{
    enhanced,breakable,
    colback=ImuPaleTeal,colframe=ImuTeal,
    boxrule=0.8pt,arc=2pt,
    left=8pt,right=8pt,top=7pt,bottom=7pt,
    title={How this chart was made},
    fonttitle=\bfseries\color{ImuNavy},
    #1
}

\title{%
    \vspace*{1.2cm}
    {\color{ImuNavy}\Huge\bfseries IMU Error Model\\[0.25em]
    \color{ImuBlue}\Large Example Profiles and Analysis}\par
    \vspace{1.0cm}
    \color{ImuSlate}\large Visual examples for the mathematical specification\par
    \vfill
    \color{ImuSlate}\normalsize Version 0.1.2\\
    Analysis based on the example profile set\\
    \today
}
\author{}
\date{}

\begin{document}
    \maketitle
    \tableofcontents


    \chapter{Purpose and scope}

    This companion document gives the example profiles a visual, comparative
    context. It shows how the same \software{ImuModel} interface can represent
    consumer, tactical, navigation, industrial, and biological reference cases,
    then traces those profiles through Allan-deviation analysis, sampled
    measurements, thermal sweeps, and unaided dead reckoning.

    The plots show the behavior produced by the example profiles. They provide a
    consistent basis for visual inspection and comparison; they are not vendor
    performance claims, certification evidence, or measurements of physical
    hardware.

    \begin{notionalbox}
        Every hardware-oriented profile in this document is a notional, approximate
        modeling estimate derived from public datasheets or other cited source material.
        Its structured source metadata records the evidence used to choose values, but
        the resulting profile is not authoritative, an official manufacturer
        configuration, certification evidence, or a guarantee of instrument
        performance. The HumanVestibular profile is a comparison benchmark rather than
        hardware.
    \end{notionalbox}


    \section{Reproducible artifact set}

    The figures in this document are PNG assets under
    \path{artifacts/showcase/}. They can be regenerated with:

    \begin{verbatim}
python3 scripts/ci.py showcase
    \end{verbatim}

    This command runs seeded simulations through the public loading and model
    interfaces, computes Allan deviation or reconstruction errors, and saves each
    figure as both PNG and SVG. The PNG files are used here for portable PDF
    generation; the matching SVG files remain available for vector inspection.


    \chapter{The example profile set}

    The repository contains fifteen canonical YAML profile documents. They all use
    the same validated \software{ProfileDocument} envelope: model name, nominal
    sample period, structured provenance, and accelerometer/gyroscope channel
    configuration. The classes below describe the intended modeling category, not
    a certification tier.

    \begin{longtable}{@{}L{1.35in}L{1.18in}L{0.92in}C{0.72in}p{1.95in}@{}}
        \caption{Example profiles included in the repository.}\label{tab:profiles}\\
        \toprule
        \textbf{Profile} & \textbf{Family} & \textbf{Class} & \textbf{$\Delta t$ (s)} & \textbf{Role in the examples} \\
        \midrule
        \endfirsthead
        \toprule
        \textbf{Profile} & \textbf{Family} & \textbf{Class} & \textbf{$\Delta t$ (s)} & \textbf{Role in the examples} \\
        \midrule
        \endhead
        \bottomrule
        \endfoot
        ADIS16470 & ADIS16470 & industrial & 0.0005 & Datasheet-oriented inertial reference. \\
        ICM-42688-P & ICM-42688-P & consumer & 0.0010 & Consumer six-axis reference. \\
        HG1700AG58 & HG1700 & tactical & 0.0100 & Tactical reference for measurement and reconstruction. \\
        HG1700AG59 & HG1700 & tactical & 0.001667 & Tactical variant with faster sampling. \\
        HG1700AG60 & HG1700 & tactical & 0.001667 & Tactical variant for profile comparison. \\
        HG1700AG61 & HG1700 & tactical & 0.001667 & Tactical variant for profile comparison. \\
        HG1700AG71 & HG1700 & tactical & 0.0100 & Tactical variant with slower sampling. \\
        HG5700AA01 & HG5700 & navigation & 0.001667 & Navigation-grade family example. \\
        HG5700BA01 & HG5700 & navigation & 0.001667 & Navigation-grade family example. \\
        HG5700CA01 & HG5700 & navigation & 0.001667 & Representative navigation curve. \\
        HG9900 & HG9900 & navigation & 0.003333 & Navigation curve and reconstruction case. \\
        HumanVestibular & human vestibular & biological & 0.0100 & Human benchmark. \\
        iPhoneLike & phone-like & consumer & 0.0100 & Consumer curve and reconstruction case. \\
        SBG\_PULSE\_40 & PULSE-40 & tactical & 0.0005 & Thermal trend example and tactical comparison. \\
    \end{longtable}


    \section{What a profile controls}

    Each channel can combine independent or correlated white noise, turn-on bias,
    stationary Gauss--Markov or random-walk bias, finite-band flicker bias,
    misalignment, scale factor, nonlinearity, clipping, quantization, and thermal
    coefficients. Values can be scalar or explicitly per-axis where the profile
    evidence supports that distinction.

    The model receives truth velocity and orientation and returns body-frame
    velocity and rotation-vector increments. The figures therefore show the
    behavior of a sensor-error model, not a complete navigation filter. In
    particular, the reconstruction plot is a deliberately unaided dead-reckoning
    demonstration.


    \chapter{Allan-deviation views}


    \section{Profile-class comparison}

    \begin{figure}[p]
        \centering
        \includegraphics[width=0.98\textwidth]{hero-allan-ladder.png}
        \caption{Representative accelerometer and gyroscope Allan-deviation curves
        across consumer, tactical, navigation, and phone-like example profiles.}
        \label{fig:allan-ladder}
    \end{figure}

    \begin{methodbox}
        For each selected profile, the script loads the canonical YAML document,
        constructs an \software{ImuModel} with a deterministic random seed, samples
        zero-motion rates, and computes overlapping block-mean differences over a
        logarithmic set of averaging times. The curves compare model behavior under a
        common analysis interface; they do not imply that every source supplied a full
        Allan-deviation characterization.
    \end{methodbox}

    Figure~\ref{fig:allan-ladder} is the broadest visual summary in the artifact
    set. The vertical placement reflects configured error magnitudes, while the
    shape reflects the composition of white noise, persistent bias processes, and
    flicker-like components. The sample period matters: it determines both the
    available shortest averaging time and the number of simulated samples in a
    fixed-duration run.


    \section{Error-process anatomy}

    \begin{figure}[p]
        \centering
        \includegraphics[width=0.92\textwidth]{error-anatomy-allan.png}
        \caption{Allan signatures of isolated white noise, Gauss--Markov bias,
            finite-band flicker, and their combined fixture.}
        \label{fig:error-anatomy}
    \end{figure}

    This fixture is intentionally inspectable rather than tied to a hardware
    profile. It shows why a single scalar noise-density summary is insufficient for
    longer averaging times: white noise, a finite-correlation bias, and a finite
    flicker bank leave different signatures. The combined curve is produced by the
    same stateful sampler used for the hardware-oriented profiles.


    \section{Reading the Allan plots}

    The plots use Allan deviation, not power spectral density. A falling short-term
    curve is consistent with white rate noise; a shoulder or plateau can arise from
    correlated bias or flicker-like terms; and a rising long-term curve indicates
    nonstationary or low-frequency error behavior. The exact interpretation depends
    on the configured process, duration, cluster schedule, and the axis selected
    for the rendering.


    \chapter{Measurements and thermal response}


    \section{Truth versus sampled measurement}

    \begin{figure}[p]
        \centering
        \includegraphics[width=0.98\textwidth]{measurement-time-series.png}
        \caption{A deterministic HG1700AG58 motion rendered as imperfect sampled
        accelerometer and gyroscope measurements, with residuals in practical units.}
        \label{fig:measurement}
    \end{figure}

    The truth trajectory is deterministic: a slowly varying acceleration and a
    rotating body frame. The measured channels include the configured profile
    effects and are plotted in the body-frame convention of the package. The right
    column expresses residuals in micro-$g$ and degrees per hour to make small
    errors visible; those display conversions do not change the stored model units.


    \section{Temperature sweep}

    \begin{figure}[p]
        \centering
        \includegraphics[width=0.98\textwidth]{thermal-trends.png}
        \caption{Thermal-only accelerometer and gyroscope trends for the notional
        SBG PULSE 40 profile, swept from $-40$ to $71$ degrees Celsius.}
        \label{fig:thermal}
    \end{figure}

    The thermal rendering disables stochastic and unrelated deterministic terms,
    uses a constant reference motion, and supplies the current temperature to the
    channel thermal model. Temperature is an explicit Celsius input at the public
    sampling boundary. The reference temperature is $25\,^{\circ}\mathrm{C}$, so a
    zero crossing near that point is expected for a purely linear coefficient.

    \begin{notionalbox}
        The thermal plot is a sensitivity visualization, not a thermal qualification
        test. It makes configured coefficients visible and does not model thermal lag,
        sensor self-heating, or environmental coupling.
    \end{notionalbox}


    \chapter{Unaided reconstruction}

    \begin{figure}[p]
        \centering
        \includegraphics[width=0.98\textwidth]{reconstruction-drift.png}
        \caption{Position and attitude error accumulated by dead reckoning the same
        truth motion through three representative error profiles.}
        \label{fig:reconstruction}
    \end{figure}

    The reconstruction experiment integrates the returned body-frame increments
    without aiding, external observations, or a Kalman filter. It compares HG9900,
    HG1700AG58, and iPhoneLike examples against the same constant-acceleration,
    constant-rate truth trajectory. Position error is reported in meters and
    attitude error in degrees.

    This figure is useful for intuition, but it should not be read as a navigation
    performance ranking. The result depends on the chosen truth motion, duration,
    seed, initialization, integration convention, and the fact that no aiding is
    available. A real estimator owns the navigation-state covariance and would
    combine the sensor process with its own dynamics and observations.


    \section{A compact interpretation}

    The five views answer different questions:

    \begin{itemize}
        \item Figure~\ref{fig:allan-ladder} compares configured stochastic behavior
        across profile classes.
        \item Figure~\ref{fig:error-anatomy} separates the signatures of the core
        error processes.
        \item Figure~\ref{fig:measurement} shows how deterministic truth becomes a
        sampled body-frame measurement.
        \item Figure~\ref{fig:thermal} isolates temperature-dependent terms.
        \item Figure~\ref{fig:reconstruction} shows how increment errors accumulate
        when no navigation correction is applied.
    \end{itemize}

    Together they provide a visual companion to the configuration and estimator
    contracts without claiming more than the code and profile metadata support.


    \chapter{Reproduction and artifact manifest}

    The showcase figures can be regenerated from the project root:

    \begin{verbatim}
python3 scripts/ci.py showcase
    \end{verbatim}

    The command writes the following paired PNG and SVG assets:

    \small
    \begin{tabularx}{0.97\textwidth}{@{}L{1.85in}Y@{}}
\toprule
\textbf{Asset} & \textbf{Description} \\
\midrule
\path{hero-allan-ladder} & Representative sensor-class Allan-deviation comparison. \\
\path{error-anatomy-allan} & Allan signatures of composable error terms. \\
\path{reconstruction-drift} & Dead-reckoning position and attitude error growth. \\
\path{measurement-time-series} & Truth rates compared with imperfect sampled rates and residuals. \\
\path{thermal-trends} & Thermal-only accelerometer and gyroscope error trends. \\
\bottomrule
    \end{tabularx}
    \normalsize

    The canonical asset descriptions are also written to
    \path{artifacts/showcase/manifest.csv}. The numerical profile definitions live
    bundled under \path{imu_error_model/data/example_profiles/hardware_estimates/}; their structured source
    URLs, archive links, review dates, and metadata should be consulted before
    using any value outside an illustrative simulation.

\end{document}
